\clearpage
\fancyhead[L]{\small VOYAGE DE HUMBOLDT ET BONPLAND.}
\hypertarget{native-517-lower}{}\label{native-517-lower}
\begin{center}\small Göttingische gelehrte Anzeigen. Stück 91. Seite 898..904. 1809 Juni 10.\\[3pt]\rule{.72\linewidth}{.3pt}\end{center}

\emph{Voyage} \textsc{de Humboldt} \emph{et} \textsc{Bonpland}. \emph{Quatrième Partie, Astronomie et Magnetisme. Premier Volume, contenant un recueil d'observations astronomiques, d'opérations trigonométriques, et de mesures barométriques, faites pendant le cours d'un voyage aux régions équinoxiales du Nouveau-Continent, depuis} 1799 \emph{jusqu'en} 1803. I. Lieferung 1808, II. Lieferung 1809. gr. Quart 279 Seiten. Bei \textsc{Schöll} und \textsc{Cotta} in Paris und Tübingen.

Die Wichtigkeit und die erstaunliche Menge der geographischen Ortsbestimmungen, welche unser berühmter Landsmann \textsc{von Humboldt} auf seiner Reise in die Amerikanischen Tropenländer gemacht hat, würden allein schon hinreichen, dieser Reise unter den gelehrten Expeditionen nach der neuen Welt einen der ersten Plätze anzuweisen: der Umstand, dass so ungemein viel \emph{von Einem Manne} geleistet ist, und wiederum, dass alles dies nur Einer von den Zwecken war, die dieser seltene Mann mit gleicher Wärme umfasst, machen die Unternehmung einzig. Nichts war daher mehr zu wünschen, als dass der Schatz von Beobachtungen, welche Hr. \textsc{von Humboldt} nach Europa zurückgebracht hat, nicht flüch\SourceLeaf{520}tig und oberflächlich, sondern mit aller möglichen Gewissenhaftigkeit und Sorgfalt discutirt und benutzt werden möchte. Glücklicherweise fand Hr. \textsc{von Humboldt} in Hrn. \textsc{Oltmanns} einen Mann, der diesem weitläuftigen und beschwerlichen Geschäfte gewachsen war und sich ihm mit musterhafter Treue und ausdauerndem Fleisse unterzogen hat. Hr. \textsc{Oltmanns} hat schon früher von diesen rühmlichen Eigenschaften schätzbare Proben abgelegt: Hr. \textsc{von Humboldt} konnte ohne Bedenken die Ableitung der Resultate aus seinen Beobachtungen diesem jungen Astronomen anvertrauen. Allein Hr. \textsc{von Humboldt} hat sich damit nicht begnügt: er hat vielmehr, damit Jedermann im Stande sei, die Güte seiner Beobachtungen, die Richtigkeit der \textsc{Oltmanns}'schen Rechnungen und die verhältnissmässige Zuverlässigkeit der Resultate selbst zu würdigen, den Entschluss gefasst, seine sämmtlichen Originalbeobachtungen in extenso neben den Hauptmomenten der \textsc{Oltmanns}'schen Rechnungen durch den Druck bekannt zu machen. Gewiss sind diese seltenen, für die Geographie Amerika's so wichtigen und unter so ungewöhnlichen Beschwerden und Gefahren errungenen Beobachtungen dieser Ehre sehr würdig, wenn auch \emph{die meisten} derselben nicht von einer solchen Natur sind, dass man in Zukunft noch merklich genauere Resultate aus ihnen ziehen zu können hoffen dürfte, als sich schon jetzt aus ihnen ziehen lassen.

Von den Instrumenten, welche Hr. \textsc{von Humboldt} auf seiner Reise mit sich führte, gibt dieser, einer Nachricht des Hrn. \textsc{Oltmanns} zufolge, in der Einleitung zu gegenwärtigem Werke selbst die nöthige Auskunft. Wir wissen nicht, ob diese Einleitung bereits gedruckt ist, und sehen also nur aus einigen gelegentlich beigebrachten Notizen, dass ein Chronometer von \textsc{Berthoud} und ein Sextant von \textsc{Ramsden} die am meisten gebrauchten Instrumente waren; zuweilen wurde von einem kleinen zweizolligen \textsc{Troughton}'schen Sextanten Gebrauch gemacht, und einige male finden wir auch eines kleinen \textsc{Bird}schen Quadranten erwähnt. Die Breitenbestimmungen gründen sich am häufigsten auf Sonnenhöhen, seltener auf Sternhöhen: die Längenbestimmungen sind grösstentheils chronometrisch, bei den wichtigern Punkten aber sind auch Jupiterstrabantenverfinsterungen und Mondsdistanzen, in Cumana auch einmal eine Sonnenfinsterniss, beobachtet. Begreiflich können also die gemessenen Breiten nicht auf den Grad von Genauigkeit Anspruch machen, den man mit festen Instrumenten oder mit Multiplicationskreisen erreicht, und unter den Längenbestimmungen wird keine sich bis auf eine Minute in Bogen ganz verbürgen lassen. Allein dies kann natürlich den Werth dieser Beobachtungen in einem Lande nicht schmälern, wo die Lage vieler der vom Hrn. \textsc{von Humboldt} bestimmten Oerter bisher auf mehrere Grade ungewiss, und bei den meisten noch niemals bestimmt war.

Da es dem Zwecke unsrer Gel. Anz. nicht angemessen sein würde, unserm Reisenden Schritt vor Schritt zu folgen, und alle seine Bestimmungen umständlich zu erzählen, so begnügen wir uns, den Inhalt der drei Bücher, aus welchen vorliegende zwei Lieferungen bestehen, summarisch anzuzeigen. Das \emph{erste Buch} begreift den Aufenthalt in Spanien und die Ueberfahrt nach Südamerika: Hr. \textsc{von Humboldt} bestimmte hier folgende Punkte: Barcellona, Montserrat, Col de Balaquet, Venta de la Sienita, Valencia, Ruinen von Sagunt, Madrid, Aranjuez, Corunna und St. Croix auf Teneriffa. Von einigen dieser Punkte sind auch bereits gute Bestimmungen durch andre Beobachter vorhanden, deren vortreffliche Uebereinstimmung mit den \textsc{von Humboldt}'schen beweiset, wie viel Vertrauen letztere verdienen. Hr. \textsc{Oltmanns} hat hier, wie in der Folge bei allen Oertern, von denen andre, ältere oder neuere, Bestimmungen da waren, diese mit den \textsc{Humboldt}'schen zusammengestellt, um sie durch diese zu bestätigen oder zu berichtigen. Besonders wichtig und schätzbar sind diese Erörterungen bei denjenigen Oertern, welche als die Cardinalpunkte von allen in Südamerika gemachten Längenbestimmungen angesehen werden müssen, nemlich Cumana und Caraccas: diese Untersuchungen machen daher einen beträchtlichen Theil des \emph{zweiten} und \emph{dritten Buches} aus. Die Resultate, welche \textsc{Oltmanns} nach einer sehr sorgfältigen Discussion der \textsc{Humboldt}'schen Beobachtungen für diese beiden wichtigen Punkte herausbringt, sind folgende:
\SourceLeaf{521}
\begin{center} \begin{tabular}{@{}lll@{}} \textit{Cumana}, Länge & $4^{\mathrm h}\,26^{\mathrm m}\,0^{\mathrm s}$ & Br. $10^\circ\,27'\,52''$ nordl. \\ \textit{Caraccas}, Länge & $4\quad 37\quad 40$ & Br. $10\quad 30\quad 50$ nordl. \\ \end{tabular} \end{center}
\noindent Den übrigen Theil des zweiten Buchs füllen die Bestimmungen der Inseln Tabago und Trinidad, die die Mündungen des Orinoko und verschiedene Oerter an der Küste von Neu-Andalusien, so wie der übrige Theil des dritten Buchs die grosse Reise ins Innere von Südamerika zum Gegenstande hat, auf welcher Hr. \textsc{von Humboldt} eine grosse Anzahl merkwürdiger Punkte am Orinoko und Cassiquiare bestimmte und die Verbindung des erstern durch letztern und den Rio Negro mit dem Amazonenflusse ausser Zweifel setzte. Dieser Theil der \textsc{Humboldt}'schen Reise ist um so wichtiger, da alle bisherigen Karten von diesen Gegenden im höchsten Grade fehlerhaft waren und zum Theil sogar die Breiten um mehr als 4 Grade unrichtig angaben.

Den Werth astronomischer Beobachtungen aus diesen Theilen von Amerika kann man erst dann gehörig würdigen, wenn man die mannigfaltigen Schwierigkeiten erwägt, mit denen dort der Beobachter zu kämpfen hat. In Caraccas ist das Wetter so unbeständig, dass man es als ein höchst seltenes Glück ansehen muss, wenn einmal eine erwartete Beobachtung gelingt. Von sieben und zwanzig Nächten, welche Hr. \textsc{von Humboldt} zu durchwachen die Geduld hatte, um Verfinsterungen von Jupiterstrabanten zu beobachten, waren die meisten ganz verloren. Zuweilen ist der Himmel 5 Minuten vor einer erwarteten Erscheinung noch ganz klar, und augenblicklich überzogen. Den 16. Dec. 1799, wo ein Trabant verfinstert werden sollte, hatte Hr. \textsc{von Humboldt} den Planeten bei völlig heiterem Himmel bereits im Felde des Fernrohrs, sein Gehülfe \textsc{Bonpland} zählte schon die Secunden, nur Eine Minute noch, und die Beobachtung wäre gemacht gewesen: allein auf einmal entsteht ein so dichter Nebel, dass von dem Planeten selbst nichts mehr zu sehen ist. In den Wüsteneien des Orinoko ist die Lage des Beobachters noch ausserordentlicher. Während der Nacht, schreibt Hr. \textsc{von Humboldt}, wurden wir ringsum von dem Gebrüll der Tiger geängstigt, welches von dem Geheul unsrer Hunde accompagnirt wurde. Die Crocodile, angelockt von dem Feuer, welches unsre Indianischen Führer am Ufer angezündet hatten, richteten sich mit halbem Leibe aus dem Wasser empor, als wollten sie uns beobachten. Aber selbst gegen diese Scenen wird man endlich durch häufige Gewohnheit gleichgültig.

Dem zweiten Buche ist noch unter dem Titel: \emph{Essai sur les refractions dans la zone torride}, eine Abhandlung des Hrn. \textsc{von Humboldt} angehängt, worin die interessante Frage, ob in den verschiedenen Climaten einerlei Refractionstafel angewandt werden dürfe? aus verschiedenen, zum Theil neuen Gesichtspunkten betrachtet wird. Hr. \textsc{von Humboldt} zieht zuerst den Einfluss des verschiedenen Mischungsverhältnisses der Bestandtheile der atmosphärischen Luft und ihren hygrometrischen Zustand in Betrachtung und zeigt, dass man keine Ursache habe, hieraus eine verschiedene Refraction für verschiedene Climate zu vermuthen. Wichtiger ist der Einfluss des Gesetzes der Wärmeabnahme in den höheren Luftschichten auf die Refractionen in kleinen Höhen und im Horizont. Die Frage, ob dieses Gesetz in den verschiedenen Erdzonen dasselbe sei, ist daher von der grössten Wichtigkeit. Hr. \textsc{von Humboldt} hat hierüber mehrere sehr merkwürdige unmittelbare Versuche angestellt, die in Ansehung jener Wärmeabnahme in der heissen Zone eine überraschend grosse Uebereinstimmung geben; das Resultat aus allen ist 122 Toisen Höhe für die Wärmeabnahme von 1 Grad Reaumur: dies ist fast gar nicht verschieden von dem, was \textsc{Gay-Lussac} in einem Luftball in einer Höhe von 6980 Meter über Paris fand. Es scheint daher, dass, in der warmen Jahrszeit, dieses Gesetz in der gemässigten Zone ganz dasselbe ist, wie in der heissen. Bei tieferem Thermometerstande hingegen scheint die Wärme nach oben zu langsamer abzunehmen: in Ermangelung directer Beobachtungen hierüber fügt Hr. \textsc{von Humboldt} diejenigen Resultate bei, welche Hr. \textsc{Mathieu} aus zwei bekannten, in Torneå bei sehr starker Kälte ange\SourceLeaf{522}stellten, Beobachtungen mit Hülfe der \textsc{Laplace}'schen Hypothese über die Abnahme der Dichtigkeit der Luft abgeleitet hat, und welche für die Höhe, wo die Wärme 1 Grad Reaumur abnimmt, $156\frac{1}{2}$ Toisen geben. Nach diesem Resultate sollte man also glauben, dass nicht die Refractionen in der heissen Zone, sondern die in der kalten, eine andere Tafel erfordern, als in der gemässigten, was gerade das Gegentheil von dem ist, was einerseits \textsc{Bouguer}, und andrerseits \textsc{Maupertuis} und \textsc{Lemonnier} behauptet haben. Indess weiss man, dass \textsc{Bouguer} selbst an seinen Beobachtungen geändert hat, um sie mit den vermeinten kleinern Refractionen in Uebereinstimmung zu bringen; ferner sind auch \textsc{Legentil}'s ebenfalls in der heissen Zone, in Pondichery, angestellte und von \textsc{Delambre} aufs neue berechnete Beobachtungen dagegen, welche sich recht gut mit \textsc{Bradley}'s, aber nicht mit \textsc{Bouguer}'s, Tafel vertragen; und endlich hat Hr. \textsc{von Humboldt} selbst mehrere Beobachtungen über die Refraction in Cumana, Caraccas, Acapulco und auf der Insel Cuba angestellt, welche durchgehends grössere Zahlen geben, als die \textsc{Bouguer}'sche Tafel. Nach allen diesen Gründen scheint es allerdings das beste zu sein, so lange, bis einmal in der heissen Zone von einem geübten Beobachter, und mit Werkzeugen, wie sie der neuern Beobachtungskunst angemessen sind, ein Cursus von zahlreichen, und die verschiedensten Zustände der Atmosphäre umfassenden, Beobachtungen angestellt sein wird, die in der gemässigten Zone construirte Refractionstafel auch in der heissen Zone beizubehalten. Dagegen wird in der gemässigten und kalten Zone der Umstand, dass die Abnahme der Wärme in den höheren Luftschichten desto langsamer ist, je tiefer das Thermometer in der untersten Luftschicht steht, die Nothwendigkeit einer Modification der Art, wie bei unsern Refractionstafeln auf die Temperatur Rücksicht genommen wird, nach sich ziehen. Leider! wird freilich die Hoffnung, hier je ins Klare zu kommen, durch die Bemerkung \textsc{Delambre}'s sehr niedergeschlagen, dass die Horizontalrefractionen bei fast ungeändertem Barometer- und Thermometerstande öfters Veränderungen von 4 Minuten leiden, ohne dass man irgend einen Grund davon anzugeben weiss.

Bemerken müssen wir noch, dass der Druck dieses \textsc{Humboldt}'schen Werkes mit einer dem innern Werthe desselben angemessenen äussern Eleganz ausgeführt ist. Wir wünschten, dies auch von der Genauigkeit des Drucks selbs in den Zahlangaben rühmen zu können: allein ungern haben wir bemerkt, dass in dieser Hinsicht die Correctur ziemlich nachlässig gemacht ist. Da der Abdruck von Originalbeobachtungen nur in so fern Werth hat, als man in alle Zahlen Vertrauen setzen kann, so ist zu erwarten, dass der Herausgeber noch einmal eine sorgfältige Vergleichung mit der Handschrift machen und alle Druckfehler am Ende gewissenhaft anzeigen werde.
